To move beyond point-estimate reporting and inform hardware procurement and mission doctrine, we couple a six-parameter global sensitivity audit with a multi-objective Pareto exploration across $\BParetoConfigs{}$ verified dispatch configurations and end-to-end offline verification inside our dark 3D digital twin command dashboard.

\subsection{Global Parameter Sensitivity and Pairwise Interactions}
We evaluate six physical and operational parameters—gimbal elevation ceiling $\beta_{\max}$, odometry drift slope $b$, communication outage duration $\Delta t_{\mathrm{out}}$, gimbal pan rate $\omega_{\mathrm{pan}}$, base position standard deviation $a$, and vehicle ground speed $v$—across symmetric $\pm 20\%$ perturbations around their nominal values (and the physically valid interval $[76^\circ,82^\circ]$ for $\beta_{\max}$). To aggregate cross-stage responsiveness onto a unified dimensionless scale, we form a normalized composite performance index $J$ weighting P1 quality, P2 response latency, P3 admissible distance, and P4 disturbance recovery at $(0.25,\,0.30,\,0.25,\,0.20)$; we note explicitly that while this weighting reflects balanced operational preference across all four mission dimensions, the top-tier parameter ranking is dictated primarily by steep physical gradients within individual subproblems rather than arbitrary weight tuning.

Table~\ref{tab:sens} and Figure~\ref{fig:sens} confirm a sharp Pareto hierarchy among the six drivers. Gimbal elevation ceiling $\beta_{\max}$ ranks first with a $\BGlobalTopSwing{}\%$ composite swing ($34.16\%$ of total main-effect share), directly reflecting the $\sec^2\beta_{\max}$ curvature derived in Eq.~\eqref{eq:p1-quality}. Odometry drift slope $b$ ranks second ($10.42\%$ swing, $26.00\%$ share) because position drift accounts for $84.61\%$ of pointing variance at $s_{\max}$, followed closely by communication outage duration $\Delta t_{\mathrm{out}}$ ($9.26\%$ swing, $23.12\%$ share). Together, these top three factors explain \textbf{83.28\%} of overall system sensitivity (Figure~\ref{fig:sens}d). By contrast, gimbal pan rate $\omega_{\mathrm{pan}}$ ($2.88\%$ swing) nearly doubles the impact of vehicle ground speed $v$ ($1.57\%$ swing), providing quantitative proof that short urban canyon legs are predominantly slew-limited rather than chassis-speed-limited. The $6\times 6$ pairwise interaction matrix (Figure~\ref{fig:sens}c, audited for exact symmetry with $0.0$ max asymmetry) remains modest across all off-diagonal pairs, confirming that first-order single-parameter upgrades deliver nearly additive engineering gains.

\ReviewFigure{fig09_global_sensitivity.pdf}{94mm}
{Global parameter sensitivity suite: (a) OAT main-effect ranking, (b) tornado response chart, (c) $6\times 6$ pairwise interaction matrix, and (d) Pareto share with cumulative contribution.}
{fig:sens}

\subsection{Multi-Objective Pareto Tradeoff and 3D Digital Twin Verification}
Dispatch response latency $t_{\mathrm{last}}$ is not the sole operational concern: excessive fleet driving distance accelerates odometry drift (Problem 3) and battery consumption. Sweeping shortlist depths $k\in\{2,3,4,5,6,8,10,12\}$ and beam widths $w\in\{1,5,10,20,40\}$ across $\BParetoConfigs{}$ verified solver configurations identifies $\BParetoFrontCount{}$ non-dominated Pareto solutions in $(t_{\mathrm{last}},\,\text{total fleet travel})$ space (Figure~\ref{fig:pareto3d}a). Examining this frontier reveals a nuanced, transparent engineering tradeoff between two configurations that tie at the exact discrete optimum $t_{\mathrm{last}}=\BBeamEight{}$ s: setting $(k=5,w=20)$ achieves slightly lower aggregate fleet travel ($36.15$ m, maximum single-vehicle excursion $10.44$ m) and therefore sits directly on the two-objective Pareto boundary, whereas our nominal $(k=8,w=20)$ configuration travels $38.24$ m (max single-vehicle excursion $12.53$ m, still $<19\%$ of $s_{\max}=\BIntegratedCrossing{}$ m) while achieving closer mean observation range ($59.92$ m vs.\ $62.51$ m) and higher average geometric quality ($0.02008$ m\textsuperscript{$-1$} vs.\ $0.01937$ m\textsuperscript{$-1$}). At the ultra-low-travel extreme, the Nearest baseline trims total travel to $21.11$ m at the expense of $+0.690$ s slower interception ($\BNearestEight{}$ s). Table~\ref{tab:deploy} translates these verified tradeoffs into four concrete operational deployment profiles.

\ReviewFigure{fig10_pareto_and_3d.pdf}{96mm}
{(a) Multi-objective Pareto frontier across \BParetoConfigs{} verified solver configurations (\BParetoFrontCount{} non-dominated); (b) dark 3D digital twin command dashboard keyframes verified offline in Chromium.}
{fig:pareto3d}

Finally, to eliminate hidden geometric or scheduling bugs across all 74 interactive timelines (nominal plus \BCaseCount{} disturbance cases), we built a fully self-contained, zero-CDN dark-theme 3D digital twin dashboard (\texttt{B\_Simulator.html}) embedding Plotly WebGL 3D alongside an automatic isometric Canvas 2.5D fallback and a one-click \emph{Auto-Tour: V5 Fault \& Dual-Strategy Demo} mode. Automated offline verification in Headless Chromium (with network access disabled and both WebGL and WebGL-disabled Canvas modes exercised, including forced WebGL context loss) recorded zero external requests, zero page errors, 770/770 passing state assertions, and \textbf{0 building collisions across all 892 executed travel segments and 4,503 stationary states}. Figure~\ref{fig:pareto3d}b displays the three verified 3D keyframes—normal coordination ($3.093$ s), V5 fault trigger ($5.000$ s), and reschedule recovery ($12.927$ s)—framed for crisp monochrome and color print legibility.
